% Acknowledgments draft. The first paragraph is the personal one; the rest
% are the acknowledgment texts the data providers ask for. Fill the [...]
% placeholders and confirm the affiliations and roles before submission.

\begin{acknowledgments}
We are grateful to Prof.\ Zheng Zheng for his guidance throughout this
project, for many discussions that shaped the analysis, and in particular
for proposing the subset test that established the sample-variance origin
of the outer-profile offset in the simulated single-galaxy stack. We thank
Dr.\ Di He for [helpful discussions on / a careful reading of an early draft
of / advice on] this work. % CHECK: state Dr. He's contribution specifically
Any remaining errors are our own.

Funding for SDSS-III has been provided by the Alfred P.\ Sloan Foundation,
the Participating Institutions, the National Science Foundation, and the
U.S.\ Department of Energy Office of Science. The SDSS-III web site is
\url{http://www.sdss3.org/}. SDSS-III is managed by the Astrophysical
Research Consortium for the Participating Institutions of the SDSS-III
Collaboration. % Paste the full institution list required by SDSS-III here.

This work is based on observations obtained with Planck
(\url{http://www.esa.int/Planck}), an ESA science mission with instruments
and contributions directly funded by ESA Member States, NASA, and Canada.
We use the Planck 2018 lensing products from the Planck Legacy Archive.

The CosmoSim database used in this paper is a service by the
Leibniz-Institute for Astrophysics Potsdam (AIP). The MultiDark database
was developed in cooperation with the Spanish MultiDark Consolider Project
CSD2009-00064. The authors gratefully acknowledge the Gauss Centre for
Supercomputing e.V.\ (\url{www.gauss-centre.eu}) and the Partnership for
Advanced Supercomputing in Europe (PRACE, \url{www.prace-ri.eu}) for
funding the MultiDark simulation project by providing computing time on the
GCS Supercomputer SuperMUC at Leibniz Supercomputing Centre
(LRZ, \url{www.lrz.de}).

This research made use of \textsc{astropy}, \textsc{healpy}, \textsc{numpy},
\textsc{scipy}, and \textsc{matplotlib}. % add \citep keys if your journal wants software cited
\emph{Declaration of AI use.} Generative AI tools were used in this project.
This declaration states what they were used for, what they were not used for,
and when.

\emph{Tools and versions.} All AI assistance was obtained through Claude Code,
Anthropic's command-line coding assistant. The models used, as recorded in the
co-authorship trailers of the project's version-control history, were Claude
Opus 4.6 (February 2026), Claude Opus 5, including its 1M-context configuration
(August and September 2026), and Claude Fable 5.1 (September 2026).
% CHECK: add any other AI tool used outside this repository. If none, say so.

\emph{Work completed without AI assistance.} The measurement method of this
paper was developed before any AI tool was used on the project. Between
2025 July 17 and 2026 February 5, the first author wrote the original analysis
in Jupyter notebooks and committed it in 59 revisions: the construction of the
convergence maps, the galaxy-pair catalogs, the stacking and symmetrization
procedure, the control subtraction, and the first error estimates. Bugs found
and fixed in that period, including an error in the Mpc$/h$ conversion and a
comparison of the equatorial and Galactic coordinate implementations, were
diagnosed by the first author. The research question, the choice to look for
filaments between CMASS pairs in the Planck lensing map, and the design of the
control were proposed and refined in discussion with Prof.\ Zheng Zheng and
Dr.\ Di He. These notebooks are preserved in the project repository under
\texttt{legacy/}.

\emph{Work in which AI assisted.} From 2026 February onward, AI assistance was
used at three stages.

First, software. The notebooks were restructured into the current command-line
pipeline with AI assistance, and much of that code---catalog loading, pair
selection, the stacking routines, the jackknife covariance estimation, and the
plotting scripts---was corrected with AI assistance. Debugging was a continuous
part of this: corrections made with AI assistance include an out-of-memory failure
in the random-catalog loader, a normalization error in the control-pair map, and a
half-pixel error in the centering of the stacking grid.

Second, verification. AI was used to audit the code and the intermediate
results. Two errors that affected numbers previously written into this draft
were found this way: a mismatch between the transverse separation bin used for
the mock and that used for the data (Section~\ref{sec:sim}), and an
inconsistency between the bridge-band definitions applied to the mock and to
the observation (Section~\ref{sec:method}). Both were corrected before
submission.

Third, the manuscript. Sections of this report were polished with AI assistance
from the authors' outlines, results, and notes, and were then revised more by the
authors. AI was also used for a spelling-convention pass.

\emph{What AI was not used for.} No numerical result reported here was produced
by a language model. Every value in the text, tables, and figures is an output
of the analysis code running on the public BOSS, Planck, and MultiDark data
described in Section~\ref{sec:data}, and can be recomputed from the archived
pipeline. AI was not used to generate, alter, or select data; the catalogs and
maps are public and unmodified. The physical model was not chosen by AI: the
cosmology is Planck 2018 and the halo occupation form is that of
\citet{White2011}, adopted on the advice of Prof.\ Zheng Zheng. No AI-edited
code or text entered the repository or this report without being read and
accepted by the authors, who are responsible for all of its content, including
any remaining errors.

\emph{Timing and frequency.} AI assistance was used between 2026 February and
2026 September, in two concentrated periods: the restructuring of the pipeline
in 2026 February, and the analysis, verification, and writing in 2026 August
and September. No AI tool was used on this project before 2026 February.

\end{acknowledgments}
